% GENERATED FILE. Edit canonical lessons, then run npm run generate:books.
% canonical-source-hash: fnv1a64:f0bd3f73c04a2e20
% canonical-lessons: AR-C457-takhfid, AR-C457-qasima, AR-C457-khaffada, AR-C457-rakhusa, AR-C457-anfaqa

\chapter{The sales}
\label{ch:ar-the-sales}
\hlchaptermodality{457}

\begin{chapteropening}
\textbf{By the end of this chapter:} \emph{I can say a price cut, a voucher, to lower, to get cheaper and to spend money.}

Complete the last lesson of chapter 457: I can say a price cut, a voucher, to lower, to get cheaper and to spend money.
\end{chapteropening}

\section[\texorpdfstring{\ar{تخفيض} (takhfīḍ)}{تخفيض (takhfīḍ)}]{\ar{تخفيض} (takhfīḍ) — a price cut; a reduction}

\label{lesson:AR-C457-takhfid}



\begin{quote}
Before the new one: say the Arabic for to fit, then the Arabic for to be too tight.
\end{quote}

\subsection*{You'll want to know: \ar{تخفيض}}
\textbf{\ar{تخفيض}} — \emph{takhfīḍ} — \textquotedblleft{}a price cut; a reduction\textquotedblright{}.

\begin{quote}
\ar{اشتريت} \ar{المعطف} \ar{بتخفيض} \ar{كبير}.
\end{quote}

\begin{quote}
\emph{ishtaraytu al-miʿṭaf bi-takhfīḍ kabīr.}
\end{quote}

\begin{quote}
I bought the coat with a big price cut.
\end{quote}

The plural \emph{takhfīḍāt} is the sales.

\begin{cousinweb}
\textbf{\ar{تخفيض}} (\emph{takhfīḍ}) is built on the native Arabic root \emph{kh-f-ḍ}, \textquotedblleft{}to be low; to lower\textquotedblright{} (Lane). It is the verbal noun of Form II.
\end{cousinweb}

\begin{grammarlens}[title={What you've built}]
One more word for getting around.
\end{grammarlens}

\subsection*{Your turn}
\begin{itemize}
\raggedright
  \item \emph{Say it:} \emph{takhfīḍ}
  \item \emph{Say it:} \emph{takhfīḍ}, once more
  \item \emph{Say it:} \emph{ishtaraytu al-miʿṭaf bi-takhfīḍ kabīr.} — I bought the coat with a big price cut.
  \item \emph{Recall:} say the Arabic for to fit, then the Arabic for to be too tight, then say \emph{takhfīḍ} again
\end{itemize}

\begin{tcolorbox}[breakable,colback=teal!4,colframe=teal!35!black,title={Before you move on}]
What does \ar{تخفيض} mean? (\textquotedblleft{}A price cut; a reduction\textquotedblright{}.) Say it once more.
\end{tcolorbox}
\section[\texorpdfstring{\ar{قسيمة} (qasīma)}{قسيمة (qasīma)}]{\ar{قسيمة} (qasīma) — a voucher; a coupon}

\label{lesson:AR-C457-qasima}



\begin{quote}
Before the new one: say the Arabic for to be too tight, then the Arabic for a price cut.
\end{quote}

\subsection*{You'll want to know: \ar{قسيمة}}
\textbf{\ar{قسيمة}} — \emph{qasīma} — \textquotedblleft{}a voucher; a coupon\textquotedblright{}.

\begin{quote}
\ar{استعملت} \ar{أمي} \ar{قسيمة} \ar{في} \ar{المتجر}.
\end{quote}

\begin{quote}
\emph{istaʿmalat ummī qasīma fī al-matjar.}
\end{quote}

\begin{quote}
My mother used a voucher in the shop.
\end{quote}

\begin{cousinweb}
\textbf{\ar{قسيمة}} (\emph{qasīma}) is built on the native Arabic root \emph{q-s-m}, \textquotedblleft{}to divide\textquotedblright{} (Lane). It shares its root with \emph{qism}, \textquotedblleft{}a portion, a share\textquotedblright{} (Lane); the sense of a voucher or a coupon is modern.
\end{cousinweb}

\begin{grammarlens}[title={What you've built}]
One more word for getting around.
\end{grammarlens}

\subsection*{Your turn}
\begin{itemize}
\raggedright
  \item \emph{Say it:} \emph{qasīma}
  \item \emph{Say it:} \emph{qasīma}, once more
  \item \emph{Say it:} \emph{istaʿmalat ummī qasīma fī al-matjar.} — My mother used a voucher in the shop.
  \item \emph{Recall:} say the Arabic for to be too tight, then the Arabic for a price cut, then say \emph{qasīma} again
\end{itemize}

\begin{tcolorbox}[breakable,colback=teal!4,colframe=teal!35!black,title={Before you move on}]
What does \ar{قسيمة} mean? (\textquotedblleft{}A voucher; a coupon\textquotedblright{}.) Say it once more.
\end{tcolorbox}
\section[\texorpdfstring{\ar{خفض} (khaffaḍa)}{خفض (khaffaḍa)}]{\ar{خفض} (khaffaḍa) — to lower; to reduce}

\label{lesson:AR-C457-khaffada}



\begin{quote}
Before the new one: say the Arabic for a price cut, then the Arabic for a voucher.
\end{quote}

\subsection*{You'll want to know: \ar{خفض}}
\textbf{\ar{خفض}} — \emph{khaffaḍa} — \textquotedblleft{}to lower; to reduce\textquotedblright{}.

\begin{quote}
\ar{خفض} \ar{المتجر} \ar{ثمن} \ar{الأحذية}.
\end{quote}

\begin{quote}
\emph{khaffaḍa al-matjar thaman al-aḥdhiya.}
\end{quote}

\begin{quote}
The shop lowered the price of shoes.
\end{quote}

Present: \ar{يخفض}, \emph{yukhaffiḍ}.

\begin{cousinweb}
\textbf{\ar{خفض}} (\emph{khaffaḍa}) is built on the native Arabic root \emph{kh-f-ḍ}, \textquotedblleft{}to be low; to lower\textquotedblright{} (Lane). It is Form II; \emph{munkhafiḍ}, \textquotedblleft{}low\textquotedblright{}, shares the root.
\end{cousinweb}

\begin{grammarlens}[title={What you've built}]
One more word for getting around.
\end{grammarlens}

\subsection*{Your turn}
\begin{itemize}
\raggedright
  \item \emph{Say it:} \emph{khaffaḍa}
  \item \emph{Say it:} \emph{khaffaḍa}, once more
  \item \emph{Say it:} \emph{khaffaḍa al-matjar thaman al-aḥdhiya.} — The shop lowered the price of shoes.
  \item \emph{Recall:} say the Arabic for a price cut, then the Arabic for a voucher, then say \emph{khaffaḍa} again
\end{itemize}

\begin{tcolorbox}[breakable,colback=teal!4,colframe=teal!35!black,title={Before you move on}]
What does \ar{خفض} mean? (\textquotedblleft{}To lower; to reduce\textquotedblright{}.) Say it once more.
\end{tcolorbox}
\section[\texorpdfstring{\ar{رخص} (rakhuṣa)}{رخص (rakhuṣa)}]{\ar{رخص} (rakhuṣa) — to get cheaper}

\label{lesson:AR-C457-rakhusa}



\begin{quote}
Before the new one: say the Arabic for a voucher, then the Arabic for to lower.
\end{quote}

\subsection*{You'll want to know: \ar{رخص}}
\textbf{\ar{رخص}} — \emph{rakhuṣa} — \textquotedblleft{}to get cheaper\textquotedblright{}.

\begin{quote}
\ar{رخصت} \ar{الفاكهة} \ar{في} \ar{الصيف}.
\end{quote}

\begin{quote}
\emph{rakhuṣat al-fākiha fī aṣ-ṣayf.}
\end{quote}

\begin{quote}
Fruit got cheaper in the summer.
\end{quote}

Present: \ar{يرخص}, \emph{yarkhuṣ}.

\begin{cousinweb}
\textbf{\ar{رخص}} (\emph{rakhuṣa}) is built on the native Arabic root \emph{r-kh-ṣ}, \textquotedblleft{}to be soft, easy, cheap\textquotedblright{} (Lane). It is the plain verb beside \emph{rakhīṣ}, \textquotedblleft{}cheap\textquotedblright{}.
\end{cousinweb}

\begin{grammarlens}[title={What you've built}]
One more word for getting around.
\end{grammarlens}

\subsection*{Your turn}
\begin{itemize}
\raggedright
  \item \emph{Say it:} \emph{rakhuṣa}
  \item \emph{Say it:} \emph{rakhuṣa}, once more
  \item \emph{Say it:} \emph{rakhuṣat al-fākiha fī aṣ-ṣayf.} — Fruit got cheaper in the summer.
  \item \emph{Recall:} say the Arabic for a voucher, then the Arabic for to lower, then say \emph{rakhuṣa} again
\end{itemize}

\begin{tcolorbox}[breakable,colback=teal!4,colframe=teal!35!black,title={Before you move on}]
What does \ar{رخص} mean? (\textquotedblleft{}To get cheaper\textquotedblright{}.) Say it once more.
\end{tcolorbox}
\section[\texorpdfstring{\ar{أنفق} (anfaqa)}{أنفق (anfaqa)}]{\ar{أنفق} (anfaqa) — to spend money}

\label{lesson:AR-C457-anfaqa}



\begin{quote}
Before the new one: say the Arabic for to lower, then the Arabic for to get cheaper.
\end{quote}

\subsection*{You'll want to know: \ar{أنفق}}
\textbf{\ar{أنفق}} — \emph{anfaqa} — \textquotedblleft{}to spend money\textquotedblright{}.

\begin{quote}
\ar{أنفقنا} \ar{كل} \ar{مالنا} \ar{في} \ar{السوق}.
\end{quote}

\begin{quote}
\emph{anfaqnā kulla mālinā fī as-sūq.}
\end{quote}

\begin{quote}
We spent all our money at the market.
\end{quote}

Present: \ar{ينفق}, \emph{yunfiq}.

\begin{cousinweb}
\textbf{\ar{أنفق}} (\emph{anfaqa}) is built on the native Arabic root \emph{n-f-q}, \textquotedblleft{}to pass away, to run out\textquotedblright{} (Lane). It is Form IV, \textquotedblleft{}to make money run out\textquotedblright{}; Lane gives \emph{nafaq}, \textquotedblleft{}a tunnel\textquotedblright{}, on the same root.
\end{cousinweb}

\begin{grammarlens}[title={What you've built}]
One more word for getting around.
\end{grammarlens}

\subsection*{Your turn}
\begin{itemize}
\raggedright
  \item \emph{Say it:} \emph{anfaqa}
  \item \emph{Say it:} \emph{anfaqa}, once more
  \item \emph{Say it:} \emph{anfaqnā kulla mālinā fī as-sūq.} — We spent all our money at the market.
  \item \emph{Recall:} say the Arabic for to lower, then the Arabic for to get cheaper, then say \emph{anfaqa} again
\end{itemize}

\begin{tcolorbox}[breakable,colback=teal!4,colframe=teal!35!black,title={Before you move on}]
What does \ar{أنفق} mean? (\textquotedblleft{}To spend money\textquotedblright{}.) Say it once more.
\end{tcolorbox}
